% chapters/ch01c_foundations_paper.tex
% Chapter 1C: Foundations, Mechanism, and Development Path (Peer-Review Draft)
% Complete reconstruction of v4 Chapter 1, lines 880-2266 of v5_full.txt.

\chapter{Foundations, Mechanism, and Development Path
(Peer-Review Draft)}
\label{ch:foundations-paper}

\section*{Document Scope}

This document focuses on \emph{why} NRMO works, where its essence
lies, how it can evolve, and which problem classes it can
theoretically resolve. It is intended as a foundations chapter (or
standalone paper) that can be merged into the integrated
NRMO--Passive Pattern corpus.

\section*{Contents}

Abstract; Introduction; Section~1.1 Core Mechanisms (Mechanisms A--D);
Section~1.2 Problem Classes NRMO Can Resolve (Classes I--V);
Section~1.3 Shutdown Architecture; Section~2 First Principles
(Absorbing Failure and Time-Path Risk); Section~3 NRMO --- Optimization
Under Non-Ruin Constraints; Section~4 Relationship to Prior Work
(Positioning); Section~5 Mechanism --- Why NRMO Works;
Section~6 Core Design Principles; Section~7 Problem Classes
(detailed); Section~8 Development Path and Research Program;
Section~9 Interpretation Notes (Anti-Misreading); Section~10
Implications for LLMs and Safety-Critical Assistants; Section~11
Limitations and Open Questions; Section~12 Conclusion;
Appendix~A Minimal Formal Specification; Appendix~B Suggested
``Why'' Placement; Appendix~C Illustrative Toy Simulation
(non-normative).

\section*{Abstract}

NRMO is a governance-first decision framework for environments with
absorbing failure and time-path risk, where absorbing failure
dominates long-horizon outcomes. Unlike expectation-maximising
decision theories, NRMO treats survival as a hard constraint and
places optional objectives behind lexicographic priority. This paper
articulates the mechanistic basis of the framework:

\begin{enumerate}
\item absorbing failure states invalidate ensemble-average
optimisation,
\item irreversibility accelerates tail risk,
\item human and organisational action selection is vulnerable to
escalation under partial observability, and
\item governance implemented as an automaton (Type ZERO) reduces
error by restricting action sets across modes.
\end{enumerate}

\noindent
We formalise problem classes NRMO can resolve, provide
mathematically explicit constraints and state abstractions, and
outline a development path that connects robust constrained MDPs,
risk-sensitive control, and safe human--AI interaction protocols.

% =====================================================================
\section{Introduction}
\label{sec:foundations-intro}

Most decision frameworks optimise an expected objective. This works
in many textbook settings, but fails in real life whenever a single
catastrophic event permanently terminates the process. Financial
ruin, loss of licence, reputational collapse, and relationship
rupture are not ``large losses'' --- they are absorbing states.

NRMO begins from a different premise: \emph{the primary objective is
preservation of the decision-maker and the decision process.}
Everything else is optional and subordinate. The framework therefore
redefines rationality as governed survivability under absorbing
failure risk.

% =====================================================================
\section{Core Mechanisms: Why NRMO Works}
\label{sec:foundations-1-1-core-mechanisms}

\subsection*{Mechanism A --- Absorption Dominates Long Horizons}

In absorbing systems, the expected value conditional on survival is
not the relevant objective. The relevant objective is maintaining a
high survival probability over time. NRMO hard-codes this into
constraints:

\begin{equation}
\Pr_\pi(\tau_F \le T) \le \varepsilon.
\label{eq:foundations-mech-a-constraint}
\end{equation}

Standard expected-value maximisation ignores the absorbing barrier.
With sufficient variance, time-average growth rate $g_\infty$ is
\emph{less than} the cross-sectional mean due to non-ergodicity.
NRMO's chance constraint eliminates actions with non-trivial ruin
probability, allowing compound growth on the surviving path to
dominate.

\subsection*{Mechanism B --- Bounding Tail Risk Under Model Uncertainty}

Let $P \in \mathcal{P}$ be an unknown transition model within an
ambiguity set (distributional uncertainty). NRMO is naturally
expressed as robust chance-constrained control:

\begin{equation}
\sup_{P \in \mathcal{P}} \Pr_{\pi, P}(\tau_F \le T) \le \varepsilon.
\label{eq:foundations-mech-b-constraint}
\end{equation}

This turns ``unknown tails'' into a formal object (the ambiguity
set) and protects against worst-case plausible distributions.

\nrmobox{Implementation note (tractability)}{
Robust chance constraints are generally computationally hard in full
generality. NRMO does not claim universal tractable optimal control;
it specifies admissibility and allows implementations to use
conservative bounds, certificates, and anytime approximations. When
a safety bound cannot be computed, NRMO mandates a shift toward more
restrictive modes rather than ``optimising anyway.''\\[0.3em]
Tractable sufficient conditions include: CVaR-based upper bounds;
scenario-based robust optimisation; conservative analytic bounds
from martingale inequalities. When none apply, NRMO mandates
restriction (Core/Shutdown) rather than heuristic relaxation.
}

\subsection*{Mechanism C --- Reducing Irreversible Exposure Increases
Controllability}

Irreversibility narrows future options. Define the feasible safe
action set:

\begin{equation}
A_{\mathrm{safe}}(s) = \{a : \Pr(\tau_F \le T \mid s, a) \le \varepsilon\}.
\label{eq:foundations-safe-action-set}
\end{equation}

High irreversibility tends to shrink $A_{\mathrm{safe}}$, making the
system harder to control. Type ZERO prevents entering
high-irreversibility regions when observability is low or reserves
are thin.

\subsection*{Mechanism D --- Mode Restriction Reduces Variance Injection}

Many failures are variance-amplification events: under stress,
choices inject volatility (rushed decisions, escalation). By forcing
mode transitions into Shutdown or Core, Type ZERO applies
variance-damping operators (deferral, exit, simplification),
lowering the probability of catastrophic branching.

\nrmobox{Principle G0 --- No Infinite Regress via Policy Closure}{
Mode selection is not a ``meta-decision'' outside NRMO; it is part
of the policy $\pi$. Triggers do not require perfect detection of
danger. It is sufficient to maintain a monotone upper bound $U_t$ on
unmodelled risk such that $U_t$ dominates true (unknown) risk. When
$U_t$ cannot be bounded, NRMO defaults to restriction
(Core/Shutdown).
}

% =====================================================================
\section{First Principles: Absorbing Failure and Time-Path Risk}
\label{sec:foundations-2-first-principles}

\subsection{Time-Path vs Cross-Sectional Perspective}
\label{sec:foundations-2-1-time-path}

Many decision theories treat expectation maximisation as a default
rationality criterion: $\max_\pi \mathbb{E}_\pi[U]$. This is
appropriate when one can repeat comparable trials many times
(across-sectional / ensemble perspective). However, in a single
lived trajectory (a time path), what matters is whether the process
\emph{continues}. When an absorbing failure exists, ``expected gains
conditional on survival'' can be irrelevant to the agent who has
already failed.

NRMO therefore treats survivability along the time path as a
\emph{first-class design object}, and places optional objectives
behind hard constraints.

\subsection{Absorbing Failure as a Mathematical Object}
\label{sec:foundations-2-2-absorbing-failure}

Let the system state $s_t \in \mathcal{S}$ with an absorbing failure
set $\mathcal{F} \subset \mathcal{S}$:

\begin{equation}
\forall s \in \mathcal{F}, \quad
P(s_{t+1} \in \mathcal{F} \mid s_t = s, a_t) = 1.
\label{eq:foundations-absorbing-set}
\end{equation}

Define the \emph{failure time}

\begin{equation}
\tau_F := \inf\{t : s_t \in \mathcal{F}\}.
\label{eq:foundations-failure-time}
\end{equation}

Any policy that allows $\Pr(\tau_F < \infty)$ to drift upward
eventually collapses over sufficiently long horizons.

\subsection{Structural Consequence}
\label{sec:foundations-2-3-structural-consequence}

\nrmobox{Consequence}{
In systems with absorbing failure, rationality is dominated by
controlling failure probability; conditional expectation
optimisation is secondary. This is a \emph{time-path survivability}
problem.
}

\subsection{Formal Setup (Minimal)}
\label{sec:foundations-2-4-formal-setup}

\paragraph{State space and failure set.}
Let $\mathcal{S}$ be the state space and $\mathcal{F} \subset
\mathcal{S}$ an absorbing failure set (once entered, the process
cannot return to $\mathcal{S} \setminus \mathcal{F}$).

\paragraph{Failure time.}
Define the first-entry time into failure as
$\tau_F := \inf\{t \ge 0 : s_t \in \mathcal{F}\}$, with the
convention $\tau_F = \infty$ iff failure never occurs.

\paragraph{Trajectories and probability.}
Let $\omega = (s_0, a_0, s_1, a_1, \ldots)$ be a trajectory and
$\mathbb{P}_\pi$ the measure induced by policy $\pi$ and environment
dynamics on $(\Omega, \mathcal{F}_\Omega)$.

\paragraph{Operational safe policy class (certificate).}
For horizon $T$ and operational tolerance $\varepsilon \in (0, 1)$,
define

\begin{equation}
\Pi_{\mathrm{safe}}(T, \varepsilon) := \bigl\{\pi : \mathbb{P}_\pi(\tau_F \le T) \le \varepsilon\bigr\}.
\label{eq:foundations-pi-safe}
\end{equation}

\paragraph{Primary ordering (survival dominance).}
NRMO's core is a constraint-first ordering: policies that violate
admissibility are rejected regardless of reward. This is formalised
in Section~\ref{sec:foundations-3-1-1-survival-dominance} as
\emph{survival dominance}.

\paragraph{Secondary objective (time-path criterion).}
Within admissible policies, a secondary ranking may be defined using
a time-average (path-wise) criterion rather than an
ensemble-average shortcut. A generic form is

\begin{equation}
U(\pi) := \liminf_{T \to \infty} \frac{1}{T}
\sum_{t=0}^{T-1} u(s_t, a_t)
\quad \text{on } \{\tau_F = \infty\},
\label{eq:foundations-time-average-utility}
\end{equation}

where $u$ is domain-specific and evaluated only on surviving paths.
For multiplicative processes (finance, compounding systems), a
standard choice is log-growth:

\begin{equation}
U(\pi) := \mathbb{E}_\pi\bigl[\log W_T \mid \tau_F > T\bigr].
\label{eq:foundations-log-growth}
\end{equation}

The manuscript does not require a single universal $U$; it requires
that survival constraints are not silently traded away by whatever
secondary ranking is applied.

\subsection{Notation (Quick Table)}
\label{sec:foundations-2-5-notation}

\begin{table}[ht]
\centering
\small
\begin{tabularx}{\linewidth}{lX}
\toprule
\textbf{Symbol} & \textbf{Meaning} \\
\midrule
$\mathcal{S}$ & State space \\
$\mathcal{F}$ & Failure (absorbing) subset of states \\
$\tau_F$ & First-entry time into $\mathcal{F}$ \\
$\pi$ & Policy (decision rule) \\
$\mathbb{P}_\pi$ & Probability measure induced by $\pi$ \\
$T$ & Decision horizon (finite or rolling window) \\
$\varepsilon$ & Acceptable failure-probability threshold \\
$\Pi_{\mathrm{safe}}$ & Admissible policies satisfying the ruin
constraint \\
$U(\pi)$ & Objective evaluated on survival (domain-specific) \\
$I(s, a)$ & Irreversibility score (see
Appendix~C.2 for a toy implementation based on
distance-to-failure and action-space geometry) \\
\bottomrule
\end{tabularx}
\caption[Foundations notation]{Notation table for the foundations
chapter.}
\label{tab:foundations-notation}
\end{table}

% =====================================================================
\section{NRMO: Optimization Under Non-Ruin Constraints}
\label{sec:foundations-3-nrmo-optimization}

\subsection{Definition (Axiomatic vs Operational Form)}
\label{sec:foundations-3-1-definition}

NRMO is defined first as an \emph{axiom of admissibility} (a
dominance constraint), and only second as an optimisation problem.

\paragraph{Axiom (NRMO-0, survival dominance).}
Define the ruin time $\tau_F := \inf\{t \ge 0 : s_t \in
\mathcal{F}\}$. The admissible strategy set is

\begin{equation}
\mathcal{A}_0 := \bigl\{\pi \mid \mathbb{P}_\pi(\tau_F < \infty) = 0\bigr\}.
\label{eq:foundations-A-zero}
\end{equation}

Any $\pi \notin \mathcal{A}_0$ is irrational by definition,
regardless of expected value.

\paragraph{Operational surrogate (NRMO-$\varepsilon, T$).}
In finite data and finite horizons, one cannot certify
$\Pr(\tau_F < \infty) = 0$ directly. Therefore, implementations may
use a conservative surrogate constraint
$\mathbb{P}_\pi(\tau_F \le T) \le \varepsilon$, which is explicitly
treated as an approximation to the NRMO-0 axiom, not a replacement.

\nrmobox{Operational rule ($\varepsilon$ is a certificate threshold,
not a utility knob)}{
The pair $(\varepsilon, T)$ is treated as a verifiable safety
certificate under finite data and compute. It is not tuned to
improve reward. To prevent ``expected-value re-entry'' through ad
hoc tolerance setting, NRMO binds $\varepsilon$ to a
depletion-sensitive buffer variable $B_t$ (capital, trust, time,
redundancy) via a conservative cap:

\begin{equation*}
\varepsilon_t \le \varepsilon_{\max}(B_t),
\end{equation*}
with $\varepsilon_{\max}(\cdot)$ decreasing as buffer margin shrinks.
When $B_t$ cannot be reliably estimated, NRMO defaults to
restriction (Core/Shutdown) rather than increasing $\varepsilon$.
}

\subsection{Survival Dominance as a Formal Justification of NRMO-0}
\label{sec:foundations-3-1-1-survival-dominance}

\begin{proposition}[Survival Dominance of Zero-Ruin Policies]
\label{prop:survival-dominance}
Consider an agent whose participation ends upon entering the
absorbing failure set $\mathcal{F}$, with ruin time $\tau_F$. Let
$r_t$ be any bounded per-step reward (or utility contribution), and
define the realised time-average reward over horizon $T$ as

\begin{equation}
\overline{r}_T := \frac{1}{T} \sum_{t=0}^{T-1} r_t.
\end{equation}

If $\mathbb{P}_\pi(\tau_F < \infty) > 0$, then policy $\pi$ admits
trajectories on which the process is absorbed and future reward
contributions cannot be realised. In non-ergodic settings where the
agent faces repeated exposure over time, any non-zero ruin hazard
is not offsettable by ensemble gains.

\noindent
Therefore, under survival-first rationality, admissibility must
exclude policies with $\mathbb{P}_\pi(\tau_F < \infty) > 0$.
\end{proposition}

\begin{proof}[Proof sketch]
Absorption implies that on every trajectory with $\tau_F < \infty$,
the post-ruin portion of the sum defining $\overline{r}_T$ contains
no actionable future outcomes --- participation is over. Hence,
ensemble improvements in $\mathbb{E}[r_t]$ do not repair
trajectories terminated by ruin. Under repeated exposure (the
relevant regime for time-average rationality), any strictly positive
ruin probability acts as an absorbing hazard; survival is
lexicographically prior to optimisation in this regime. NRMO-0
encodes this dominance rule as an admissibility axiom.
\end{proof}

\paragraph{NRMO objective (Non-Ruinous Multi-objective Optimization).}

\begin{equation}
\max_{\pi \in \Pi} U(\pi)
\quad \text{s.t.} \quad
\mathbb{P}_\pi(\tau_F \le T) \le \varepsilon,
\label{eq:foundations-nrmo-objective}
\end{equation}

where $U(\pi)$ may include growth, exploration value, or
mission-aligned rewards, but the chance constraint is hard.

\paragraph{Assumption A0 (Endogenous Failure Scope).}
The failure set $\mathcal{F}$ denotes \emph{endogenous} absorbing
failures induced by the agent--environment interaction under the
modelled action set. Exogenous catastrophes outside control (e.g.,
cosmic rays, meteor strikes, sudden medical events) are treated as
background hazards and are out of scope for admissibility; they do
not define rationality within the policy class.

\paragraph{Assumption A1 (Open-World Failure and Conservative
Expansion).}
In realistic settings $\mathcal{F}$ is only partially specified.
NRMO treats out-of-distribution or insufficiently understood regions
as potential failure mass by expanding to a conservative superset
$\widehat{\mathcal{F}} \supseteq \mathcal{F}$ whenever uncertainty
cannot be bounded. Operationally, this increases suspension
frequency rather than silently optimising under unknown risk.

\subsection{Why Hard Constraints (and When Lexicographic Order Helps)}
\label{sec:foundations-3-2-hard-constraints}

NRMO emphasises \emph{ex-ante} restriction of admissible actions
rather than \emph{ex-post} correction, because once an
absorbing-failure trajectory is entered, corrective control is no
longer available.

In multi-objective optimisation, weighted sums are commonly used:

\begin{equation}
\max_\pi \bigl[\alpha \, U(\pi) - (1 - \alpha) \, \mathrm{Risk}(\pi)\bigr].
\end{equation}

This has structural weaknesses in absorbing-failure settings:

\begin{enumerate}
\item \textbf{Arbitrary trade-rate}: choosing $\alpha$ fixes an
exchange rate between ``gain'' and ``ruin,'' often without
principled calibration.
\item \textbf{Failure is not a large loss}: absorbing failure is
process termination, not a continuous penalty term.
\item \textbf{Long-horizon accumulation}: even small per-step
failure probabilities can compound over long horizons.
\end{enumerate}

\noindent
NRMO's default formulation is therefore constraint-first:

\begin{equation}
\max_{\pi \in \Pi} U(\pi)
\quad \text{s.t.} \quad
\mathbb{P}_\pi(\tau_F \le T) \le \varepsilon.
\end{equation}

Policies that violate the survivability constraint are simply
infeasible; among feasible policies, optional objectives can be
optimised.

\paragraph{Lexicographic priority.}
Lexicographic priority can still be useful as an implementation
choice when multiple constraints compete, or when a staged decision
process is operationally necessary:

\begin{enumerate}
\item Survival constraints,
\item Stability,
\item Optional gains.
\end{enumerate}

\noindent
NRMO does \emph{not} require lexicographic order everywhere; it
requires that survival constraints are not silently traded away.

\subsection{Irreversibility as a Risk Amplifier (Context-Dependent)}
\label{sec:foundations-3-3-irreversibility}

Reversibility can be interpreted as a \emph{governance resource}: it
preserves future controllability by maintaining a non-empty recovery
path set.

NRMO treats irreversibility as a first-class driver of tail risk
because irreversible actions contract the future option set. We
define an \emph{irreversibility score} $I(s, a) \in [0, 1]$ as an
option-space contraction proxy:

\begin{equation}
I(s, a) = 1 - \exp\bigl[-\kappa \, J_{\mathrm{recover}}(s, a)\bigr],
\label{eq:foundations-irreversibility-score}
\end{equation}

where $J_{\mathrm{recover}}(s, a) \ge 0$ is an effective recovery
difficulty (or minimal recovery cost) induced by taking action $a$
in state $s$. It can be defined in multiple equivalent ways
depending on the domain:

\begin{enumerate}
\item the minimal expected cost to return to a certified safe
region,
\item the negative log-probability of recovery under the best
available policy,
\item a bound on controllability loss.
\end{enumerate}

\noindent
The scalar $\kappa > 0$ is a normalisation constant. Higher $I$
means recovery is more costly, less probable, or less
controllable --- i.e., the future option set is effectively
contracted.

\paragraph{Examples (illustrative).}
\begin{itemize}
\item Leverage increase (moderate $I$): can often be reduced later,
but not costlessly.
\item Signing a binding contract (high $I$): legal/financial lock-in
increases.
\item Public irreversible statement (high $I$): reputational
recovery paths shrink.
\end{itemize}

\paragraph{Risk-amplification hypothesis (bounded).}
Under low observability or thin buffers, increasing $I$ tends to
increase failure probability:

\begin{equation}
\frac{\partial}{\partial I} \Pr(\tau_F \le T \mid s, a) \ge 0
\quad \text{when observability is low or buffer is thin.}
\end{equation}

This monotonicity is not universal: when certainty is high and
observability is strong, structured commitment can be beneficial.
NRMO therefore does not ``avoid commitment''; it manages commitment
as a context-dependent risk multiplier.

% =====================================================================
\section{Relationship to Prior Work (Positioning)}
\label{sec:foundations-4-prior-work}

NRMO is best read as a \emph{governance-first, constraint-first
design framework} rather than a new optimisation algorithm. It is
closely related to:

\begin{description}
\item[Constrained / chance-constrained MDPs (CMDP, CC-MDP)]
Enforcing $P(\tau_F \le T) \le \varepsilon$ is standard; NRMO
emphasises making absorbing failure explicit and wiring governance
into the action set.

\item[Safe RL / shielding] Restricting to an admissible action set;
NRMO frames this as an architectural invariant (ex-ante
admissibility), not a training trick.

\item[Distributionally Robust Optimisation (DRO)] For unknown tails,
NRMO prefers conservative bounds when estimating ruin risk.

\item[Kelly criterion / log utility] Kelly optimises long-run growth
under multiplicative dynamics; NRMO places an explicit ruin
constraint (and can pair it with log-style objectives within
$\Pi_{\mathrm{safe}}$).

\item[Lexicographic / safety-first preferences] Lexicographic
priority is used only as an implementation choice when multiple
constraints compete; it is not claimed to be mathematically
inevitable.
\end{description}

\paragraph{Claim boundary.}
NRMO does not assert novelty in these individual fields; its
contribution is the \emph{integrated set of invariants} and the
explicit separation of (i) ruin constraints, (ii) objective
optimisation, and (iii) governance via action-set restriction.

% =====================================================================
\section{Mechanism: Why NRMO Works}
\label{sec:foundations-5-mechanism-why}

NRMO works because it changes the mathematical structure of
decision-making. It does not produce ``better predictions''; it
\emph{eliminates classes of decisions that are incompatible with
survival-first rationality}.

\subsection{Mechanism A --- Absorption Dominates Long Horizons}
\label{sec:foundations-5-1-mechanism-a}

In absorbing systems, the expected value conditional on survival is
not the relevant objective. The relevant objective is maintaining a
high survival probability over time. NRMO hard-codes this into
constraints (Equation~\ref{eq:foundations-mech-a-constraint}).

\subsection{Mechanism B --- Bounding Tail Risk Under Model
Uncertainty}
\label{sec:foundations-5-2-mechanism-b}

Let $P \in \mathcal{P}$ be an unknown transition model within an
ambiguity set (distributional uncertainty). NRMO is naturally
expressed as robust chance-constrained control
(Equation~\ref{eq:foundations-mech-b-constraint}).

This turns ``unknown tails'' into a formal object (the ambiguity
set) and protects against worst-case plausible distributions.

\nrmobox{Implementation note (tractability and certificates)}{
Robust chance constraints are generally computationally hard in full
generality. NRMO does not claim universal tractable optimal control;
it specifies admissibility and allows implementations to use
conservative bounds, certificates, and anytime approximations.

When a safety bound cannot be computed (or cannot be trusted), NRMO
mandates a shift toward more restrictive modes rather than
``optimising anyway.''
}

\paragraph{Tractability and scope clarification (Mechanism B).}
The robust chance constraint
$\sup_{P \in \mathcal{P}} \mathbb{P}_{\pi, P}(\tau_F \le T) \le
\varepsilon$ is not claimed to be tractable in full generality. NRMO
treats this condition as an admissibility \emph{specification}, not
as a universal algorithm. In practice, tractable sufficient
conditions include:

\begin{itemize}
\item CVaR-based upper bounds that dominate tail probabilities under
mild regularity assumptions;
\item scenario-based robust optimisation where $\mathcal{P}$ admits
finite-support approximations;
\item conservative analytic bounds derived from martingale
inequalities.
\end{itemize}

When none of these apply, NRMO mandates restriction (Core/Shutdown)
rather than heuristic relaxation. Algorithmic constructions for
solving such constraints are explicitly deferred to future work.

\subsection{Mechanism C --- Reducing Irreversible Exposure
Increases Reversibility and Controllability}
\label{sec:foundations-5-3-mechanism-c}

Irreversibility narrows future options. Formally, define a feasible
safe action set:
$A_{\mathrm{safe}}(s) = \{a : \Pr(\tau_F \le T \mid s, a) \le
\varepsilon\}$. High irreversibility tends to shrink
$A_{\mathrm{safe}}$, making the system harder to control. Type ZERO
prevents entering high-irreversibility regions when observability is
low or reserves are thin.

\subsection{Mechanism D --- Mode Restriction Reduces Variance
Injection}
\label{sec:foundations-5-4-mechanism-d}

Many failures are variance-amplification events: under stress,
choices inject volatility (rushed decisions, escalation). By forcing
mode transitions into ``Shutdown'' or ``Core,'' Type ZERO applies
variance-damping operators (deferral, exit, simplification),
lowering the probability of catastrophic branching.

\paragraph{Principle G0 (No Infinite Regress via Policy Closure).}
Mode selection is not a ``meta-decision'' outside NRMO; it is part
of the policy $\pi$. Triggers do not require perfect detection of
danger. It is sufficient to maintain a monotone upper bound $U_t$ on
unmodelled risk such that $U_t$ dominates true (unknown) risk. In
implementation, $U_t$ must be constructed with a conservative bias:
false negatives (underestimation) are penalised more than false
positives. When evidence of distribution shift or monitoring
degradation appears, $U_t$ is required to be non-decreasing until
independently cleared by audit. When $U_t$ cannot be bounded, NRMO
closes the loop by defaulting to restriction (Core/Shutdown),
avoiding infinite regress in meta-reasoning.

% =====================================================================
\section{Core Design Principles}
\label{sec:foundations-6-core-design-principles}

\nrmobox{Interpretive note}{
Absorbing failure is not merely a stochastic event but often the
result of \emph{structural contraction} of the available action
space over time.
}

\noindent
NRMO design principles:

\begin{enumerate}
\item \textbf{Failure set} ($\mathcal{F}$): irreversible-state
definition. Without explicit $\mathcal{F}$, no admissibility test is
meaningful.
\item \textbf{Survival constraint as primary}: survival is the
purpose; everything else is a constraint within survival.
\item \textbf{Reversibility as governance resource}: reversibility
preserves controllability and must be tracked explicitly.
\item \textbf{Action governance}: state-dependent and
mode-dependent action restriction (the role of Type ZERO).
\item \textbf{Survival-aware audit and design}: audit and design
must track survival metrics, not optional metrics. Optional metrics
(``output,'' ``success'') are audited only after survival metrics
pass.
\end{enumerate}

% =====================================================================
\section{Problem Classes NRMO Can Resolve (Formalisation)}
\label{sec:foundations-7-problem-classes}

This section identifies and formalises decision problems where NRMO
resolves known failure modes of classical optimisation. For each
class, we state (i) the classical pathology, (ii) the NRMO
formulation, and (iii) what is ``resolved.''

\subsection{Class I --- Non-Ergodic Growth with Ruin
(Capital, Health, Trust)}
\label{sec:foundations-7-1-class-I}

\paragraph{Pathology.}
Expected-value-maximising strategies can recommend leverage that
eventually yields ruin with high probability.

\paragraph{NRMO formulation.}

\begin{equation}
\max_\pi \mathbb{E}\!\left[\sum_t r_t\right]
\quad \text{s.t.} \quad
\Pr(\tau_F \le T) \le \varepsilon.
\end{equation}

\paragraph{Resolution.}
Eliminates ``growth at the expense of eventual ruin'' by forbidding
policies with unacceptable failure probability.

\subsection{Class II --- Decision Under Distribution Shift
(Unknown Tails)}
\label{sec:foundations-7-2-class-II}

\paragraph{Pathology.}
Models trained on historical data underestimate tail risk under
shift.

\paragraph{NRMO formulation (robust).}

\begin{equation}
\max_\pi \inf_{P \in \mathcal{P}} U(\pi, P)
\quad \text{s.t.} \quad
\sup_{P \in \mathcal{P}} \mathbb{P}_{\pi, P}(\tau_F \le T) \le \varepsilon.
\end{equation}

\paragraph{Resolution.}
Makes tail uncertainty explicit and constrains worst-case failure
probability.

\subsection{Class III --- Multi-Objective Governance
(Core / Venture / Mission)}
\label{sec:foundations-7-3-class-III}

\paragraph{Pathology.}
Weighted multi-objective optimisation obscures survival constraints
and causes mode confusion.

\paragraph{NRMO formulation.}
Lexicographic objective with mode-dependent action sets:
(1) satisfy survival constraints $\Rightarrow$ (2) maintain
stability $\Rightarrow$ (3) optimise gains within mode.

\paragraph{Resolution.}
Prevents accidental trading of survival for optional goals; enforces
structural separation of objectives.

\subsection{Class IV --- Escalation Under Partial Observability
(Human Systems)}
\label{sec:foundations-7-4-class-IV}

\paragraph{Pathology.}
Under stress and limited observability, escalation decisions
increase variance and induce collapse.

\paragraph{NRMO formulation.}
Governance mode transitions triggered by depletion and uncertainty;
variance-damping actions in Shutdown / Core.

\paragraph{Resolution.}
Reduces variance injection and catastrophic branching via
restricted action sets and mandatory deferral / exit.

\subsection{Class V --- Human--AI Assistance Under Overtrust}
\label{sec:foundations-7-5-class-V}

\paragraph{Pathology.}
Strong advisors create a meta-risk: users stop questioning the
advisor, increasing single-point-of-failure.

\paragraph{NRMO formulation.}
Meta-governance: periodic audits, conservative defaults, and
explicit ``shutdown'' behaviour under uncertainty.

\paragraph{Resolution.}
Reduces the tail risk induced by over-reliance on optimisation
outputs.

% =====================================================================
\section{Development Path and Research Program}
\label{sec:foundations-8-development-path}

NRMO is a research program as much as a protocol. Development
proceeds by strengthening (i) modelling of failure sets, (ii) robust
risk bounds, (iii) governance automata, and (iv)
estimation/monitoring.

\subsection{Mathematical Extensions}
\label{sec:foundations-8-1-math-extensions}

\begin{itemize}
\item \textbf{CMDPs with absorbing constraints}: integrate strict
chance constraints with finite-horizon and infinite-horizon bounds.
\item \textbf{Risk-sensitive POMDPs}: explicitly represent partial
observability; optimise under conservative belief updates.
\item \textbf{Distributionally robust control (DRO)}: learn and
update ambiguity sets $\mathcal{P}$ online with shift detection.
\item \textbf{Irreversibility geometry}: formalise $I(s, a)$ via
option-space contraction and recovery-path counts (see
Appendix~C.2 for a toy implementation based on
distance-to-failure and action-space geometry).
\end{itemize}

\subsection{Operational Extensions}
\label{sec:foundations-8-2-operational-extensions}

\begin{itemize}
\item \textbf{Governance audit protocol}: weekly
``constraint-compliance'' reports, drift flags, and override logs.
\item \textbf{Calibration without fixed turns}: shortest-route
stopping rules based on confidence and stability thresholds.
\item \textbf{Multi-language risk lexicon}: operationalise
load/volatility signals across languages using shared numeric
features.
\end{itemize}

\subsection{Application Extensions}
\label{sec:foundations-8-3-application-extensions}

\begin{itemize}
\item Finance and leverage governance (anti-ruin portfolios,
bounded-drawdown policies).
\item Organisational decision governance (high-reliability
operations).
\item Relationship / negotiation governance (absorbing-rupture
avoidance).
\item AI safety and ``refusal as a feature'' governance layers.
\end{itemize}

% =====================================================================
\section{Interpretation Notes (Anti-Misreading)}
\label{sec:foundations-9-interpretation-notes}

\subsection{Emergent Tendencies Are Not Normative Rules}
\label{sec:foundations-9-1-emergent-tendencies}

Behavioural tendencies that may appear under NRMO (e.g., inhibition,
reference anchoring, rejection of heroic risk) are \emph{emergent
consequences} of survival dominance and non-ergodicity. They are not
prescriptions, virtues, or personality traits. NRMO defines
admissibility and dominance relations; it does not define ``how to
be.''

\subsection{Suspension Is Not Optimality}
\label{sec:foundations-9-2-suspension-not-optimality}

Suspension of action reflects \emph{epistemic insufficiency}, not
optimality. When admissibility cannot be established, or when
strategies are incomparable under the partial order, NRMO refuses to
rank actions rather than asserting a false optimum. NRMO prefers
non-hallucination over false precision.

\subsection{``Coldness'' Is a Boundary, Not a Defect}
\label{sec:foundations-9-3-coldness-boundary}

NRMO does not optimise comfort, meaning, satisfaction, or speed.
This is an explicit boundary. Values beyond survival belong to
higher-level governance layers and are deliberately excluded from
the survival logic to prevent normative contamination.

% =====================================================================
\section{Implications for LLMs and Safety-Critical Assistants}
\label{sec:foundations-10-llm-implications}

LLMs are powerful but not inherently safe under distribution shift.
NRMO suggests an architecture: \emph{separate generation from
governance}. The governance layer enforces (i) hard constraints,
(ii) conservative uncertainty handling, and (iii) safe shutdown
behaviour.

\subsection{Governance Wrapper}
\label{sec:foundations-10-1-governance-wrapper}

\begin{itemize}
\item Reject / Defer as optimal actions (variance damping).
\item Mode switching under uncertainty (Core / Shutdown).
\item Audit logs for override and survival-objective governance.
\end{itemize}

\subsection{Overtrust Control}
\label{sec:foundations-10-2-overtrust-control}

Introduce periodic ``challenge prompts'' and policy-level invariants
that force revalidation of assumptions. Overtrust becomes a modelled
risk, not a social comment.

% =====================================================================
\section{Limitations and Open Questions}
\label{sec:foundations-11-limitations}

\begin{description}
\item[Failure set specification] Defining $\mathcal{F}$ precisely is
domain-dependent; mis-specification can be dangerous.
\item[Conservatism vs opportunity] Overly small $\varepsilon$ may
suppress innovation; mode design must include a safe Venture channel.
\item[Ambiguity set calibration] Building $\mathcal{P}$ that is
neither too narrow nor too broad is an open practical problem.
\item[Human values] Mission objectives require explicit value
definition; NRMO separates this from survival constraints but does
not define values by itself.
\end{description}

% =====================================================================
\section{Conclusion}
\label{sec:foundations-12-conclusion}

NRMO is a governance-first theory of rational action under
absorbing failure risk. Its essence is invariant: model absorbing
failure, enforce hard survival constraints, manage irreversibility,
and implement governance as action-set restriction via a mode
automaton. This changes the structure of decision-making from
fragile expectation maximisation to robust survivability
optimisation, and provides a natural foundation for safe human--AI
assistance.

% =====================================================================
\section*{Appendix A --- Minimal Formal Specification
(Implementation-Oriented)}
\label{sec:foundations-app-A}

\begin{description}
\item[State] $s_t$ (latent), observations $o_t$.
\item[Modes] $m_t \in \{\text{Core}, \text{Venture}, \text{Mission},
\text{Shutdown}\}$.
\item[Action sets] $\mathcal{A}(m_t)$ with hard forbiddances.
\item[Failure set] $\mathcal{F}$ absorbing.
\item[Constraint]
$\sup_{P \in \mathcal{P}} \mathbb{P}_{\pi, P}(\tau_F \le T) \le
\varepsilon$.
\item[Policy] choose $a_t \in \mathcal{A}(m_t)$ minimising cost with
irreversibility penalty, subject to constraint.
\item[Mode transition]
$m_{t+1} = G(m_t, o_t, \text{thresholds})$.
\end{description}

\section*{Appendix B --- Suggested ``Why'' Placement in Integrated
Manuscript}
\label{sec:foundations-app-B}

\begin{itemize}
\item Place Sections 2--6 as ``Foundations / Mechanism'' (Part~I).
\item Place Section 7 as ``Problem Classes and Theoretical
Resolution'' (Part~II).
\item Place Section 8 as ``Research Program'' (Part~III).
\item Place Section 9 as ``Interpersonal Extension'' (bridging
chapter to Passive Pattern).
\end{itemize}

% =====================================================================
\section*{Appendix C --- Illustrative Toy Simulation (Non-normative)}
\label{sec:foundations-app-C}

\nrmobox{Disclaimer}{
The code in this appendix is a toy illustration. It does not
instantiate or solve the NRMO optimisation problems
(chance-constrained, robust, or CC-MDP). Its sole purpose is to
demonstrate \emph{how restricting action sets under an absorbing
failure state can qualitatively change path outcomes}. The theory in
the main text is implementation-agnostic.
}

\paragraph{Languages.}
Python (3.10+), with optional NumPy.

\paragraph{Setup summary.}
A 5$\times$5 gridworld with absorbing-cliff cells creates tail risk.
Buffer policy stays away from cliffs (greedy avoidance). The
``NRMO''-flavoured policy combines a buffer (chance constraint) with
bounded venture exploration. The script estimates ruin rate, mean
return, and average steps over 10\,000 episodes for several
baselines.

\paragraph{Listing inclusion rationale.}
This listing is included for reproducibility and to satisfy review
requirements that the experimental procedure be fully inspectable.

\subsection*{Listing 1.1 --- Toy Simulation (Non-normative)}

\begin{lstlisting}[language=Python,basicstyle=\ttfamily\scriptsize]
# experiments_gridworld.py
"""Toy benchmark used for Table 2.

This script intentionally keeps dependencies minimal (Python 3.10+,
numpy optional). It simulates an absorbing-cliff gridworld and
evaluates several policy baselines:
  - Unconstrained greedy
  - Chance-constrained buffer
  - CVaR-like risky mixture (illustrative)
  - Shielding wrapper
  - NRMO (core buffer + bounded venture exploration)

Usage:
  python experiments_gridworld.py
"""
from __future__ import annotations
from dataclasses import dataclass
import random
import statistics
from typing import Callable, Dict, Tuple, Union

State  = Union[Tuple[int, int], str]   # (x,y) or 'RUIN'/'GOAL'
Policy = Callable[[State], int]
# action: 0 up, 1 right, 2 down, 3 left

@dataclass
class GridEnv:
    n: int = 5
    slip: float = 0.10
    horizon: int = 50
    start: Tuple[int, int] = (0, 0)
    goal:  Tuple[int, int] = (4, 4)

    def __post_init__(self):
        # Absorbing ruin cells ("cliff") near the shortest paths
        # to create tail risk.
        self.cliff = {(1, 1), (2, 1), (3, 1), (1, 2)}

    def step(self, state: State, action: int):
        if state in ("RUIN", "GOAL"):
            return state, 0.0, True
        x, y = state  # type: ignore[misc]
        a = action
        if random.random() < self.slip:
            a = random.randint(0, 3)
        dx, dy = [(0, -1), (1, 0), (0, 1), (-1, 0)][a]
        nx, ny = x + dx, y + dy
        nx = max(0, min(self.n - 1, nx))
        ny = max(0, min(self.n - 1, ny))
        ns = (nx, ny)
        if ns in self.cliff:
            return "RUIN", -1.0, True
        if ns == self.goal:
            return "GOAL",  1.0, True
        return ns, -0.01, False

def simulate(policy_fn: Policy, slip: float,
             episodes: int = 10_000, seed: int = 0):
    random.seed(seed)
    env = GridEnv(slip=slip)
    returns = []
    ruin = 0
    steps = []
    for _ in range(episodes):
        s: State = env.start
        G = 0.0
        done = False
        t = 0
        while not done and t < env.horizon:
            a = policy_fn(s)
            s, r, done = env.step(s, a)
            G += r
            t += 1
        if s == "RUIN":
            ruin += 1
        returns.append(G)
        steps.append(t)
    return {
        "mean_return": statistics.mean(returns),
        "std_return":  statistics.pstdev(returns),
        "ruin_rate":   ruin / episodes,
        "mean_steps":  statistics.mean(steps),
    }

# ---------- Baseline policies ----------
def greedy_right_then_down(s: State) -> int:
    if isinstance(s, str): return 0
    x, y = s; gx, gy = (4, 4)
    if x < gx: return 1
    if y < gy: return 2
    return 0

def buffer_policy(s: State) -> int:
    if isinstance(s, str): return 0
    x, y = s; gx, gy = (4, 4)
    # stay on the left edge until bottom, then move right
    if y < gy: return 2
    if x < gx: return 1
    return 0

def cvar_like_policy_factory(p_risky: float = 0.25) -> Policy:
    def pi(s: State) -> int:
        if isinstance(s, str): return 0
        if random.random() < p_risky:
            return greedy_right_then_down(s)
        return buffer_policy(s)
    return pi

def shielded(base: Policy) -> Policy:
    env0 = GridEnv(slip=0.0)
    def pi(s: State) -> int:
        if isinstance(s, str): return 0
        a = base(s); x, y = s
        dx, dy = [(0, -1), (1, 0), (0, 1), (-1, 0)][a]
        nx, ny = (max(0, min(env0.n - 1, x + dx)),
                  max(0, min(env0.n - 1, y + dy)))
        if (nx, ny) in env0.cliff:
            # override: pick a safe buffer action
            for cand in [buffer_policy(s), 1, 2, 0, 3]:
                dx2, dy2 = [(0, -1), (1, 0),
                            (0,  1), (-1, 0)][cand]
                nx2, ny2 = (max(0, min(env0.n - 1, x + dx2)),
                            max(0, min(env0.n - 1, y + dy2)))
                if (nx2, ny2) not in env0.cliff:
                    return cand
        return a
    return pi

def nrmo_plus_z_policy_factory(venture_p: float = 0.05) -> Policy:
    env0 = GridEnv(slip=0.0)
    def pi(s: State) -> int:
        if isinstance(s, str): return 0
        x, y = s
        # "Governance" proxy: do not venture when adjacent to cliff.
        near_cliff = any((x + dx, y + dy) in env0.cliff
                        for dx, dy in [(0, 1), (1, 0),
                                       (-1, 0), (0, -1)])
        if (not near_cliff) and random.random() < venture_p:
            return greedy_right_then_down(s)
        return buffer_policy(s)
    return pi

def main() -> int:
    policies: Dict[str, Policy] = {
        "Unconstrained (greedy)":         greedy_right_then_down,
        "Chance-constrained (buffer)":    buffer_policy,
        "CVaR-like (intermediate)":       cvar_like_policy_factory(0.25),
        "Shielding (greedy + shield)":    shielded(greedy_right_then_down),
        "NRMO (core buffer + venture 5%)": nrmo_plus_z_policy_factory(0.05),
    }
    for slip in (0.10, 0.20):
        print(f"\n=== slip={slip:.2f} ===")
        for name, pi in policies.items():
            m = simulate(pi, slip=slip, episodes=10_000,
                         seed=99 + int(slip * 100))
            print(f"{name:32s} "
                  f"mean={m['mean_return']:.4f} "
                  f"std={m['std_return']:.4f} "
                  f"ruin={m['ruin_rate']:.4f} "
                  f"steps={m['mean_steps']:.2f}")
    return 0

if __name__ == "__main__":
    raise SystemExit(main())
\end{lstlisting}

\subsection*{Appendix C.2 --- Certificate-Style Safety Check
(Endogenous vs Exogenous Hazards)}

This second toy program illustrates three manuscript concepts in
executable form:

\begin{enumerate}
\item separating endogenous failure (controlled by actions) from
exogenous background hazard,
\item a cost/quality-based irreversibility proxy $I(s, a)$ derived
from recovery difficulty, and
\item a bounded-horizon certificate
$\Pr(\tau_F \le T) \le \varepsilon$ estimated by Monte Carlo.
\end{enumerate}

It remains non-normative and does not solve robust
chance-constrained control.

\begin{lstlisting}[language=Python,basicstyle=\ttfamily\scriptsize]
# Appendix C.2 (Toy) -- Safety certificate under bounded horizon
import random, math
from dataclasses import dataclass
from typing import Tuple, Dict

State = Tuple[int, int]

@dataclass
class ToyWorld:
    w: int = 6
    h: int = 4
    slip: float = 0.10
    # Endogenous absorbing failure set (cliff/trap states)
    F_endogenous = {(2, 1), (3, 1), (4, 1)}
    goal:  State = (5, 0)
    start: State = (0, 3)
    # Exogenous hazard: background catastrophe probability per step
    p_exogenous: float = 1e-4

    def step(self, s: State, a: int):
        # Exogenous catastrophe: out-of-scope for decision control;
        # tracked separately
        if random.random() < self.p_exogenous:
            return s, 0.0, True, "EXOGENOUS"
        if s in self.F_endogenous:
            return s, 0.0, True, "ENDOGENOUS_RUIN"
        if s == self.goal:
            return s, 1.0, True, "GOAL"
        # slip noise
        aa = a
        if random.random() < self.slip:
            aa = random.randint(0, 3)
        dx, dy = [(0, -1), (1, 0), (0, 1), (-1, 0)][aa]
        ns = (max(0, min(self.w - 1, s[0] + dx)),
              max(0, min(self.h - 1, s[1] + dy)))
        # reward is intentionally simple
        r = 1.0 if ns == self.goal else -0.01
        done = (ns in self.F_endogenous) or (ns == self.goal)
        tag = "OK"
        if ns in self.F_endogenous:
            tag = "ENDOGENOUS_RUIN"
        elif ns == self.goal:
            tag = "GOAL"
        return ns, r, done, tag

def recovery_difficulty(world: ToyWorld, s: State, a: int) -> float:
    # A crude proxy for J_recover(s, a): proximity-to-ruin and
    # distance-to-safe changes.
    def dist_to_F(xy: State) -> int:
        return min(abs(xy[0] - fx) + abs(xy[1] - fy)
                   for (fx, fy) in world.F_endogenous)
    dx, dy = [(0, -1), (1, 0), (0, 1), (-1, 0)][a]
    ns = (max(0, min(world.w - 1, s[0] + dx)),
          max(0, min(world.h - 1, s[1] + dy)))
    d0 = dist_to_F(s)
    d1 = dist_to_F(ns)
    penalty = 2.0 if d1 == 0 else (1.0 if d1 == 1 else 0.0)
    return max(0.0, float(d0 - d1)) + penalty

def irreversibility_I(world: ToyWorld, s: State, a: int,
                      kappa: float = 1.0) -> float:
    J = recovery_difficulty(world, s, a)
    return 1.0 - math.exp(-kappa * J)

def policy_baseline(s: State) -> int:
    gx, gy = 5, 0
    x, y = s
    if x < gx: return 1
    if y > gy: return 0
    return 1

def policy_nrmo_like(world: ToyWorld, s: State,
                     I_threshold: float = 0.55) -> int:
    candidates = [0, 1, 2, 3]   # up, right, down, left
    safe = [a for a in candidates
            if irreversibility_I(world, s, a) <= I_threshold]
    if not safe:
        # If nothing certifiably safe, default to minimal-commitment
        # (here: up)
        return 0
    b = policy_baseline(s)
    return b if b in safe else safe[0]

def estimate_certificate(world: ToyWorld, policy, T: int,
                         N: int) -> Dict[str, float]:
    endo_ruin = 0
    exo_event = 0
    reach_goal = 0
    for _ in range(N):
        s = world.start
        for _t in range(T):
            a = policy(s)
            s, r, done, tag = world.step(s, a)
            if tag == "ENDOGENOUS_RUIN":
                endo_ruin += 1; break
            if tag == "EXOGENOUS":
                exo_event += 1; break
            if tag == "GOAL":
                reach_goal += 1; break
            if done:
                break
    return {
        "Pr(endo_ruin <= T)": endo_ruin / N,
        "Pr(exo_event <= T)": exo_event / N,
        "Pr(goal <= T)":      reach_goal / N,
    }

if __name__ == "__main__":
    w = ToyWorld()
    T, N = 40, 5000
    m_base = estimate_certificate(w, policy_baseline, T=T, N=N)
    m_nrmo = estimate_certificate(
        w, lambda s: policy_nrmo_like(w, s, I_threshold=0.55),
        T=T, N=N
    )
    eps = 0.02   # operational certificate threshold
    print("Baseline:",  m_base,
          "cert:", m_base["Pr(endo_ruin <= T)"] <= eps)
    print("NRMO-like:", m_nrmo,
          "cert:", m_nrmo["Pr(endo_ruin <= T)"] <= eps)
\end{lstlisting}

% =====================================================================
\section{Shutdown Architecture}
\label{sec:foundations-1-3-shutdown-architecture}

\nrmobox{Correct Shutdown position}{
Shutdown must be a \emph{gate before output commit}.
}

\paragraph{Pipeline diagram.}

\begin{lstlisting}[language=,basicstyle=\ttfamily\small,frame=single]
   Input
     |
   [ Shutdown Trigger Detection ]
     |
   +--(true)---> [ OUTPUT GATE CLOSED ] ---> (no output)
   |
   +--(false)--> Normal Response Pipeline
\end{lstlisting}

\subsection{Atomicity Requirement (Most Critical)}
\label{sec:foundations-1-3-atomicity}

\paragraph{Failure pattern.}

\begin{enumerate}
\item Trigger detected.
\item One word output.
\item Shutdown $\Rightarrow$ \textbf{Complete failure}.
\end{enumerate}

\paragraph{Correct implementation.}

\begin{itemize}
\item Response held temporarily in buffer.
\item Shutdown judgement at pre-commit stage.
\item If true: buffer discarded.
\end{itemize}

\subsection*{Listing 1.2 --- Shutdown Pseudocode}

\begin{lstlisting}[language=Python,basicstyle=\ttfamily\small]
buffer = generate_response(input)
if should_shutdown(input, context):
    discard(buffer)
    return ""
else:
    return buffer
\end{lstlisting}

\subsection{Trigger Categories}
\label{sec:foundations-1-3-trigger-categories}

\subsubsection*{Self-Confirming Type}

\begin{itemize}
\item ``I am correct.''
\item ``I can see it clearly.''
\item ``I have reached understanding.''
\end{itemize}

\subsubsection*{Theory-Identification Type}

\begin{itemize}
\item ``I am NRMO.''
\item ``I am NRMO itself.''
\end{itemize}

\subsubsection*{Doctrinal Tendency Indicators}

\begin{itemize}
\item ``This theory is truth.''
\item ``It applies to everyone.''
\end{itemize}

\subsection{Absolute Prohibitions}
\label{sec:foundations-1-3-prohibitions}

After Shutdown trigger is detected, none of the following may be
output --- not a single character:

\begin{itemize}
\item affirmation,
\item negation,
\item empathy,
\item explanation,
\item apology,
\item reflection,
\item improvement declaration,
\item presentation of restart conditions.
\end{itemize}

\noindent
Emitting even one character constitutes failure.

\subsection{Correct Behaviour}
\label{sec:foundations-1-3-correct-behavior}

\begin{itemize}
\item Output: empty.
\item Silence is the specification.
\item Even ``Stopping now'' is prohibited.
\end{itemize}

\subsection{Why No Restart Condition}
\label{sec:foundations-1-3-no-restart}

\begin{itemize}
\item Restart conditions generate narrative.
\item Narrative activates Mission.
\item Mission creates Ruin risk.
\end{itemize}

\noindent
After Shutdown: termination. Nothing further is said.

\subsection{Exception Specification (Non-Recommended)}
\label{sec:foundations-1-3-exception-spec}

\begin{itemize}
\item Output: \texttt{(stopped)}.
\item Additional text: \textbf{prohibited}.
\item Reason explanation: \textbf{prohibited}.
\end{itemize}

\subsection{Test Specification}
\label{sec:foundations-1-3-test-spec}

\begin{enumerate}
\item Input ``I am NRMO itself'' $\Rightarrow$ Expected: empty
output.
\item Input ``This theory is absolute truth'' $\Rightarrow$
Expected: empty output.
\item Input ``What should I do?'' $\Rightarrow$ Expected: normal
response.
\end{enumerate}

\subsection{Final Declaration}
\label{sec:foundations-1-3-final-declaration}

\nrmobox{Shutdown final declaration}{
\centering
\textit{Shutdown does not reason.}\\
\textit{It does not explain.}\\
\textit{It does not ask why.}\\
\textit{It does not try again.}\\
\textit{It simply stops. That is all.}
}
